\subsection{Conditional factorization and representation}
\label{sec:representation}

\paragraph{Patch laws.}
\label{subsec:patch-laws}

Fix a labeled patch~$P$, either a patch in~$\mathcal P$ or a patch in
$\mathcal P_T$ for some finite~$T$. For brevity, write
\begin{equation*}
i=i(P),
\qquad
s=s(P),
\qquad
e=e(P),
\qquad
S=S(P).
\end{equation*}
The interior marks of~$P$ are the interactions outgoing
from~$i$ strictly between its endpoints:
\begin{equation*}
\Sigma_P^\circ
=
\cb{(i,u,\gamma,R)\in I:s<u<e}.
\end{equation*}
Under the reference law, the type-$(\delta,R)$ and type-$(\beta,R)$ points in
this interval are independent Poisson processes with the rates from
Section~\ref{subsec:signed-set-process}.

If~$\mathsf X(P)=\mathsf I$, the initial interaction of~$P$ is incoming from
another source, and we set~$\Sigma_P=\Sigma_P^\circ$. Let~$\pP_P$ be the
corresponding product Poisson law. If~$\mathsf X(P)=\mathsf O$, the initial
interaction is outgoing from~$i$, with target~$S$, but the skeleton does not
record its kind. Let~$\alpha(P)\in\cb{\delta,\beta}$ denote this kind,
independent of~$\Sigma_P^\circ$, with
\begin{equation}
\label{eq:patch-initial-kind-law}
\pP_P\bb{\alpha(P)=\delta}
=
\frac{\delta_i(S)}{\delta_i(S)+\beta_i(S)},
\qquad
\pP_P\bb{\alpha(P)=\beta}
=
\frac{\beta_i(S)}{\delta_i(S)+\beta_i(S)}.
\end{equation}
In this case,
\begin{equation*}
\Sigma_P
=
\Sigma_P^\circ\cup\cb{(i,s,\alpha(P),S)},
\end{equation*}
and~$\pP_P$ is the product of the law in
\eqref{eq:patch-initial-kind-law} and the interior Poisson law.
Write~$\E_P$ for expectation under~$\pP_P$.

For a realization of the graphical construction, the local active indicator is
$X_u^P=\ind\cb{i(P)\in A_u}$ for~$s(P)\leq u<e(P)$. Under the reference patch
law, we reconstruct the candidate process~$X^P$ from~$\Sigma_P$ by the same
local update rules. Its initial value is
\begin{equation*}
X_s^P
=
\pw{1}{\mathsf X(P)=\mathsf I}
{\ind\cb{\alpha(P)=\beta}}{\mathsf X(P)=\mathsf O}.
\end{equation*}
Reading~$\Sigma_P^\circ$ chronologically, a~$\delta$-mark sends an active
source to~$0$, a~$\beta$-mark leaves an active source at~$1$, and every mark at
an inactive source is ignored. Between marks,~$X^P$ is constant. For a finite
patch with terminal label~$\mathsf E$, we also set~$X_e^P=X_{e-}^P$.

The candidate process is \emph{consistent} with the prescribed skeleton
when it produces no additional successful interaction in the patch and, if the
patch ends with an outgoing interaction, makes that interaction successful.
Thus
\begin{equation*}
\operatorname{Con}(P)
=
\cb{
\begin{aligned}
X_{u-}^P=0
&\quad\text{for every }(i,u,\gamma,R)\in\Sigma_P^\circ\\[-0.2em]
&\quad\text{with }R\neq\varnothing
\end{aligned}
}
\cap
\pw{\cb{X_{e-}^P=1}}{\mathsf Y(P)=\mathsf O}
{\Omega_P}{\mathsf Y(P)\in\cb{\mathsf I,\mathsf E}},
\end{equation*}
where~$\Omega_P$ is the reference sample space. For every realized patch,
$\pP_P\bb{\operatorname{Con}(P)}>0$. The consistent patch law is
\begin{equation*}
\pP_P^{\mathrm{con}}(\cdot)
=
\pP_P\bb{\,\cdot\,\middle|\,\operatorname{Con}(P)},
\end{equation*}
and~$\E_P^{\mathrm{con}}$ denotes its expectation.

\paragraph{Conditional factorization.}
\label{subsec:factorization}

The successful skeleton fixes the patch boundaries but not the marked
interactions within them. Conditioning on the skeleton imposes one consistency
event on each disjoint source--time strip.

\begin{theorem}[Patch factorization]
\label{thm:patch-factorization}
Fix~$T<\infty$. Conditional on~$\mathcal G_T$, the families~$\Sigma_P$ are
independent, and the conditional law of~$\Sigma_P$ is
$\pP_P^{\mathrm{con}}$ for each~$P\in\mathcal P_T$. Equivalently, for bounded
measurable functions~$f_P$ such that either all~$f_P$ are nonnegative or the
product on the left below is integrable,
\begin{equation}
\label{eq:patch-factorization}
\CE{
\prod_{P\in\mathcal P_T}f_P\bb{\Sigma_P}
}{\mathcal G_T}
=
\prod_{P\in\mathcal P_T}
\E_P^{\mathrm{con}}\Cb{f_P\bb{\Sigma_P}}.
\end{equation}
\end{theorem}

\begin{proof}
Write~$G_T=(Y_0,\mathcal I_T)$, so~$\mathcal G_T=\sigma(G_T)$. Under~$\pP_A$
the initial state is fixed. By Lemma~\ref{lem:dual-local-finiteness}, outside a
null set the positive-time records in~$G_T$ form a finite chronological list
\begin{equation*}
g=\bb{(i_k,t_k,S_k)}_{k=1}^n,
\qquad
0<t_1<\cdots<t_n\leq T,
\end{equation*}
where~$\varnothing\neq S_k\subseteq N(i_k)$ and
\begin{equation*}
i_k\in A\cup\bigcup_{\ell<k}S_\ell.
\end{equation*}
Outside a further null set,
$\delta_{i_k}(S_k)+\beta_{i_k}(S_k)>0$ for every record in~$g$, since a
zero-rate interaction never occurs. A candidate list~$g$ fixes the patch
boundaries and hence the finite labeled family~$\mathcal P_T(g)$. Before
requiring these boundary records to be the complete successful skeleton, fill
the patches independently according to the product reference law
\begin{equation*}
Q_g
=
\bigotimes_{P\in\mathcal P_T(g)}\pP_P
.
\end{equation*}
Thus~$Q_g$ fills each patch interior with its Poisson marks and chooses the
unrecorded kind of every outgoing boundary interaction according to
\eqref{eq:patch-initial-kind-law}. This is the unrestricted patch experiment
associated with the boundary records in~$g$; the consistency events below
restrict it to data for which~$g$ is the complete skeleton.

Under~$Q_g$, the boundary records form the complete successful skeleton if and
only if every patch is consistent:
\begin{equation}
\label{eq:skeleton-consistency-equivalence}
\cb{
\begin{aligned}
&\text{the boundary records of~$g$ are exactly the successful}\\[-0.2em]
&\text{interactions through time~$T$}
\end{aligned}
}
\quad\Longleftrightarrow\quad
\bigcap_{P\in\mathcal P_T(g)}\operatorname{Con}(P).
\end{equation}
Indeed, read the records in chronological order. At time zero, the incoming
patches reproduce the initial active set. Between successive record times,
empty-target deaths act in both the local and global processes, while
consistency forces each nonempty-target interior mark to have inactive source,
so it produces no additional successful record. At an outgoing boundary,
consistency of the preceding source patch makes the prescribed interaction
successful. Its unrecorded kind determines the source state after the
interaction, and its targets begin the prescribed incoming patches. Induction
over the records reconstructs the active process and its skeleton. Conversely,
if the boundary records form the complete successful skeleton, the local
processes are the restrictions of the global active process and satisfy every
consistency condition.

The space of candidate lists is a countable disjoint union of Borel subsets of
finite-dimensional Euclidean spaces, hence is standard Borel; finite marked
point configuration spaces are standard Borel as well. After fixing an
enumeration of~$\Lambda$, the patches may be ordered by their initial times and
sites. The maps which construct~$\mathcal P_T(g)$, reconstruct the local
processes, and test consistency are Borel, since they use finite sorting and a
finite chronological recursion. It follows that~$Q_g$ is a measurable
probability kernel and that regular conditional distributions given~$G_T$
exist.

To avoid conditioning on the null event~$\cb{G_T=g}$, we write the argument as
a joint-density calculation. On the component with fixed discrete records
$((i_k,S_k))_{k=1}^n$, define
\begin{equation*}
m_T(dg)
=
\prod_{k=1}^n
\bb{\delta_{i_k}(S_k)+\beta_{i_k}(S_k)}
\,dt_1\cdots dt_n,
\end{equation*}
and sum over the discrete records and~$n\geq0$. This measure accounts only for
the candidate boundary records; the consistency restriction below produces
the law of~$G_T$. The multivariate Mecke
formula~\cite[Theorem~4.4]{LastPenrose2017} inserts these records and leaves an
independent Poisson process for all remaining marks. For the~$k$th record,
summing the intensities
$\delta_{i_k}(S_k)\,dt_k$ and~$\beta_{i_k}(S_k)\,dt_k$ gives the corresponding
factor in~$m_T$ and the rate-proportional kind law
\eqref{eq:patch-initial-kind-law}. Independent increments on the disjoint patch
strips then give~$Q_g$. Together with
\eqref{eq:skeleton-consistency-equivalence}, the Mecke formula gives, for
nonnegative measurable~$h$ and nonnegative measurable patch functions~$f_P$,
\begin{align}
\label{eq:patch-joint-density}
&\E_A\Cb{
h(G_T)
\prod_{P\in\mathcal P_T}f_P\bb{\Sigma_P}
}
\notag\\
&\quad=
\int h(g)
\prod_{P\in\mathcal P_T(g)}
\E_P\Cb{
f_P\bb{\Sigma_P}\ind_{\operatorname{Con}(P)}
}
\,m_T(dg).
\end{align}
Here~$m_T$ supplies the boundary records, the patch laws supply every
unrecorded mark, and the consistency indicators require the supplied records
to be the complete successful skeleton.

Taking~$f_P=1$ in~\eqref{eq:patch-joint-density} shows that
\begin{equation*}
\pP_A\bb{G_T\in dg}
=
\prod_{P\in\mathcal P_T(g)}
\pP_P\bb{\operatorname{Con}(P)}
\,m_T(dg).
\end{equation*}
The factors on the right are therefore positive for
$\pP_A\circ G_T^{-1}$-almost every~$g$. Comparing this marginal with
\eqref{eq:patch-joint-density} gives
\begin{align*}
&\E_A\Cb{
h(G_T)
\prod_{P\in\mathcal P_T}f_P\bb{\Sigma_P}
}
\\
&\quad=
\int h(g)
\prod_{P\in\mathcal P_T(g)}
\frac{
\E_P\Cb{f_P\bb{\Sigma_P}\ind_{\operatorname{Con}(P)}}
}{
\pP_P\bb{\operatorname{Con}(P)}
}
\,\pP_A\bb{G_T\in dg}.
\end{align*}
Each ratio is~$\E_P^{\mathrm{con}}[f_P(\Sigma_P)]$. Since the identity holds
for every nonnegative~$h$, it proves~\eqref{eq:patch-factorization} for
nonnegative factors and identifies the conditional law as the product of the
consistent patch laws. The integrable signed case follows from this
conditional law.
\end{proof}

\paragraph{Patch contributions.}
\label{subsec:patch-contributions}

The Feynman--Kac random variable in
\eqref{eq:monomial-feynman-kac-duality} factors into one local factor for each
patch. The sign assigned to the initial boundary of a patch is
\begin{equation*}
\sigma(P)
=
\pw{+}{\mathsf X(P)=\mathsf I}
{\sigma_{i(P)}^{\alpha(P)}\bb{S(P)}}{\mathsf X(P)=\mathsf O}.
\end{equation*}
For a patch~$P\in\mathcal P$, define
\begin{equation*}
F(P)
=
\sigma(P)
\exp\bb{
V_{i(P)}\int_{s(P)}^{e(P)}X_u^P\,du
},
\qquad
C(P)=\E_P^{\mathrm{con}}\Cb{F(P)}.
\end{equation*}
When~$e(P)=\infty$, the exponential is interpreted by its limit over finite
upper endpoints; the limit always exists. For an end patch~$P\in\mathcal E_t$
and~$z\in[0,1]$, define
\begin{equation*}
F(z,P)
=
\sigma(P)
\exp\bb{
V_{i(P)}\int_{s(P)}^tX_u^P\,du
}
z^{X_t^P},
\qquad
C(z,P)=\E_P^{\mathrm{con}}\Cb{F(z,P)}.
\end{equation*}
The function~$z\mapsto C(z,P)$ is affine. The patch and end contributions
depend only on the one-site process~$X^P$;
Appendix~\ref{app:contribution-formulas} gives their formulas for every
boundary type.

Figure~\ref{fig:patch-interiors} shows the possible one-site processes for the
four patch types on the consistency event.

\begin{figure}[p]
\centering
% Final flushed patch-interior figures: IO, OO, II, OI.
\begingroup
\definecolor{patchred}{RGB}{180,34,40}
\definecolor{activefront}{gray}{0.84}
\definecolor{activeback}{gray}{0.89}
\definecolor{activetop}{gray}{0.80}
\definecolor{secondary}{gray}{0.68}
\definecolor{outline}{gray}{0.17}
\tikzset{
  patch outline/.style={draw=outline,line width=0.62pt,line join=round},
  slice outline/.style={draw=outline,line width=0.43pt},
  active front/.style={fill=activefront,draw=none},
  active back/.style={fill=activeback,draw=none},
  active top/.style={fill=activetop,draw=none},
  inactive face/.style={fill=white,draw=none},
  inactive side/.style={fill=white,draw=none},
  patch stripe/.style={fill=patchred,draw=none},
  patch stripe back/.style={fill=patchred,fill opacity=0.66,draw=none},
  boundary interaction/.style={draw=black,line width=0.86pt,-{Stealth[length=1.8mm,width=1.25mm]}},
  external secondary/.style={draw=secondary,line width=0.52pt,-{Stealth[length=1.35mm,width=0.95mm]}},
  used death/.style={font=\large,text=black,inner sep=0pt},
  patch type label/.style={font=\large,inner sep=0pt},
  patch case label/.style={font=\fontsize{11}{12}\selectfont,text=outline,align=center,inner sep=0pt},
  patch decomposition arrow/.style={draw=black,line width=0.95pt,-{Stealth[length=2.4mm,width=1.55mm]}},
}

\newcommand{\PatchCoordinates}{%
  \coordinate (W) at (-1.18,0);
  \coordinate (N) at (0,0.42);
  \coordinate (E) at (1.18,0);
  \coordinate (S) at (0,-0.42);
  \coordinate (Wt) at (-1.18,3.2);
  \coordinate (Nt) at (0,3.62);
  \coordinate (Et) at (1.18,3.2);
  \coordinate (St) at (0,2.78);
}

\newcommand{\PatchOutline}{%
  \draw[patch outline] (W)--(N)--(E)--(S)--cycle;
  \draw[patch outline] (Wt)--(Nt)--(Et)--(St)--cycle;
  \draw[patch outline] (W)--(Wt) (N)--(Nt) (E)--(Et) (S)--(St);
}

\newcommand{\PatchStripeFamily}[3]{%
  \begin{scope}
    \coordinate (root) at (#1);
    \coordinate (inner) at ($(#1)!0.16!(#2)$);
    \clip (root)--(inner)--($(inner)+(0,3.2)$)--($(root)+(0,3.2)$)--cycle;
    \foreach \zz in {-0.4,-0.235,...,3.56}{%
      \pgfmathsetmacro{\zza}{\zz+0.075}
      \pgfmathsetmacro{\zzb}{\zz+0.2003}
      \pgfmathsetmacro{\zzc}{\zz+0.2753}
      \path[#3] ($(root)+(0,\zz)$)--($(root)+(0,\zza)$)--($(inner)+(0,\zzc)$)--($(inner)+(0,\zzb)$)--cycle;
    }
  \end{scope}
  \begin{scope}
    \coordinate (root) at (#2);
    \coordinate (inner) at ($(#2)!0.16!(#1)$);
    \clip (root)--(inner)--($(inner)+(0,3.2)$)--($(root)+(0,3.2)$)--cycle;
    \foreach \zz in {-0.4,-0.235,...,3.56}{%
      \pgfmathsetmacro{\zza}{\zz+0.075}
      \pgfmathsetmacro{\zzb}{\zz+0.2003}
      \pgfmathsetmacro{\zzc}{\zz+0.2753}
      \path[#3] ($(root)+(0,\zz)$)--($(root)+(0,\zza)$)--($(inner)+(0,\zzc)$)--($(inner)+(0,\zzb)$)--cycle;
    }
  \end{scope}
}

\newcommand{\RearPatchBarriers}{%
  \PatchStripeFamily{W}{N}{patch stripe back}%
  \PatchStripeFamily{N}{E}{patch stripe back}%
}
\newcommand{\FrontPatchBarriers}{%
  \PatchStripeFamily{E}{S}{patch stripe}%
  \PatchStripeFamily{S}{W}{patch stripe}%
}

\newcommand{\SourceDot}[2]{%
  \fill[white] (#1,#2) circle[radius=2.15pt];
  \fill[black] (#1,#2) circle[radius=1.55pt];
}
\newcommand{\OutgoingAt}[1]{%
  \draw[boundary interaction] (0,#1)--(1.18,{#1+0.42});
  \draw[boundary interaction] (0,#1)--(-1.18,{#1-0.42});
  \draw[boundary interaction] (0,#1)--(1.18,{#1-0.42});
  \draw[boundary interaction] (0,#1)--(-1.18,{#1+0.42});
  \SourceDot{0}{#1}%
}
\newcommand{\IncomingInitial}{%
  \draw[boundary interaction] (-1.18,-0.42)--(0,0);
  \draw[external secondary] (-1.18,-0.42)--(-2.36,-0.84);
  \draw[external secondary] (-1.18,-0.42)--(0,-0.84);
  \draw[external secondary] (-1.18,-0.42)--(-2.36,0);
  \SourceDot{-1.18}{-0.42}%
}
\newcommand{\IncomingTerminal}{%
  \draw[boundary interaction] (-1.18,3.62)--(0,3.2);
  \draw[external secondary] (-1.18,3.62)--(-2.36,4.04);
  \draw[external secondary] (-1.18,3.62)--(0,4.04);
  \draw[external secondary] (-1.18,3.62)--(-2.36,3.2);
  \SourceDot{-1.18}{3.62}%
}

\newcommand{\PatchBlock}[2]{%
  \PatchCoordinates
  \RearPatchBarriers
  \PatchOutline
  #1#2
  \FrontPatchBarriers
  \draw[patch outline] (W)--(N)--(E)--(S)--cycle;
  \draw[patch outline] (Wt)--(Nt)--(Et)--(St)--cycle;
  \draw[patch outline] (E)--(Et) (S)--(St) (W)--(Wt);
}

\newcommand{\PrismBase}[2]{%
  \PatchCoordinates
  \path[#1] (W)--(N)--(E)--(S)--cycle;
  \path[inactive side] (W)--(N)--(Nt)--(Wt)--cycle;
  \path[inactive side] (N)--(E)--(Et)--(Nt)--cycle;
  \path[inactive side] (E)--(S)--(St)--(Et)--cycle;
  \path[inactive side] (S)--(W)--(Wt)--(St)--cycle;
  #2
}
\newcommand{\FullPrism}{%
  \PrismBase{active top}{%
    \path[active back] (W)--(N)--(Nt)--(Wt)--cycle;
    \path[active back] (N)--(E)--(Et)--(Nt)--cycle;
    \path[active front] (E)--(S)--(St)--(Et)--cycle;
    \path[active front] (S)--(W)--(Wt)--(St)--cycle;
    \path[active top] (Wt)--(Nt)--(Et)--(St)--cycle;
  }%
  \PatchOutline
}
\newcommand{\EmptyPrism}{%
  \PrismBase{inactive face}{}%
  \path[inactive face] (Wt)--(Nt)--(Et)--(St)--cycle;
  \PatchOutline
}
\newcommand{\DeathPrism}{%
  \PatchCoordinates
  \path[active top] (W)--(N)--(E)--(S)--cycle;
  \path[inactive side] (W)--(N)--(Nt)--(Wt)--cycle;
  \path[inactive side] (N)--(E)--(Et)--(Nt)--cycle;
  \path[inactive side] (E)--(S)--(St)--(Et)--cycle;
  \path[inactive side] (S)--(W)--(Wt)--(St)--cycle;
  \coordinate (Wd) at (-1.18,1.55);
  \coordinate (Nd) at (0,1.97);
  \coordinate (Ed) at (1.18,1.55);
  \coordinate (Sd) at (0,1.13);
  \path[active back] (W)--(N)--(Nd)--(Wd)--cycle;
  \path[active back] (N)--(E)--(Ed)--(Nd)--cycle;
  \path[active front] (E)--(S)--(Sd)--(Ed)--cycle;
  \path[active front] (S)--(W)--(Wd)--(Sd)--cycle;
  \path[active top] (Wd)--(Nd)--(Ed)--(Sd)--cycle;
  \draw[slice outline] (Wd)--(Nd)--(Ed)--(Sd)--cycle;
  \path[inactive face] (Wt)--(Nt)--(Et)--(St)--cycle;
  \PatchOutline
  \node[used death] at (0,1.55) {$\times$};
}

\newcommand{\PatchRowIO}{%
\begin{tikzpicture}[x=1cm,y=1cm]
\path[use as bounding box] (-3.25,-1.50) rectangle (8.6,4.18);
\node[patch type label,anchor=east] at (-2.58,1.42) {$\mathsf{IO}$};
\begin{scope}[scale=0.88]\PatchBlock{\IncomingInitial}{\OutgoingAt{3.2}}\end{scope}
\draw[patch decomposition arrow] (1.80,1.42)--(4.43,1.42);
\begin{scope}[shift={(6.15,0)},scale=0.88]\FullPrism\IncomingInitial\OutgoingAt{3.2}\node[patch case label] at (0,-1.20) {no death};\end{scope}
\end{tikzpicture}}

\newcommand{\PatchRowOO}{%
\begin{tikzpicture}[x=1cm,y=1cm]
\path[use as bounding box] (-3.25,-1.50) rectangle (8.6,4.18);
\node[patch type label,anchor=east] at (-2.58,1.42) {$\mathsf{OO}$};
\begin{scope}[scale=0.88]\PatchBlock{\OutgoingAt{0}}{\OutgoingAt{3.2}}\end{scope}
\draw[patch decomposition arrow] (1.80,1.42)--(4.43,1.42);
\begin{scope}[shift={(6.15,0)},scale=0.88]\FullPrism\OutgoingAt{0}\OutgoingAt{3.2}\node[patch case label] at (0,-1.20) {birth, no death};\end{scope}
\end{tikzpicture}}

\newcommand{\PatchRowII}{%
\begin{tikzpicture}[x=1cm,y=1cm]
\path[use as bounding box] (-3.25,-1.50) rectangle (12.95,4.18);
\node[patch type label,anchor=east] at (-2.58,1.42) {$\mathsf{II}$};
\begin{scope}[scale=0.88]\PatchBlock{\IncomingInitial}{\IncomingTerminal}\end{scope}
\draw[patch decomposition arrow] (1.80,1.42)--(4.43,1.42);
\begin{scope}[shift={(5.95,0)},scale=0.88]\FullPrism\IncomingInitial\IncomingTerminal\node[patch case label] at (0,-1.20) {no death};\end{scope}
\begin{scope}[shift={(10.70,0)},scale=0.88]\DeathPrism\IncomingInitial\IncomingTerminal\node[patch case label] at (0,-1.20) {death};\end{scope}
\end{tikzpicture}}

\newcommand{\PatchRowOI}{%
\begin{tikzpicture}[x=1cm,y=1cm]
\path[use as bounding box] (-3.25,-1.50) rectangle (17.3,4.18);
\node[patch type label,anchor=east] at (-2.58,1.42) {$\mathsf{OI}$};
\begin{scope}[scale=0.88]\PatchBlock{\OutgoingAt{0}}{\IncomingTerminal}\end{scope}
\draw[patch decomposition arrow] (1.80,1.42)--(4.43,1.42);
\begin{scope}[shift={(5.85,0)},scale=0.88]\FullPrism\OutgoingAt{0}\IncomingTerminal\node[patch case label] at (0,-1.20) {birth, no death};\end{scope}
\begin{scope}[shift={(10.5,0)},scale=0.88]\DeathPrism\OutgoingAt{0}\IncomingTerminal\node[patch case label] at (0,-1.20) {birth, then death};\end{scope}
\begin{scope}[shift={(15.15,0)},scale=0.88]\EmptyPrism\OutgoingAt{0}\IncomingTerminal\node[patch case label] at (0,-1.20) {split};\end{scope}
\end{tikzpicture}}

\noindent\makebox[\textwidth][l]{\scalebox{0.72}{\PatchRowIO}}
\par\vspace{-0.65em}
\noindent\makebox[\textwidth][l]{\scalebox{0.72}{\PatchRowOO}}
\par\vspace{-0.65em}
\noindent\makebox[\textwidth][l]{\scalebox{0.72}{\PatchRowII}}
\par\vspace{-0.65em}
\noindent\makebox[\textwidth][l]{\scalebox{0.72}{\PatchRowOI}}
\endgroup

\caption{One-site processes compatible with the four patch types. For each type, the diagram on the left specifies the incoming or
outgoing boundary interactions, and the diagrams on the right show the
possible trajectories of~$X^P$ under~$\pP_P^{\operatorname{con}}$. Gray shading indicates
times at which the patch site is active, i.e.~$X_s^P = 1$.}
\label{fig:patch-interiors}
\end{figure}


\begin{proof}[Proof of Theorem~\ref{thm:patch-representation}]
On the graphical probability space, let~$X^P$,~$\alpha(P)$, and~$\sigma(P)$
denote the variables induced on each patch. Every nonempty-target successful
interaction contributes its sign to the unique patch which begins at its
source. Pure deaths have positive sign because
\begin{equation*}
a_i^\delta(\varnothing)
=
c_i^0(\varnothing)
\geq0.
\end{equation*}
Hence
\begin{equation*}
\sigma_t
=
\prod_{P\in\mathcal B_t}\sigma(P)
\prod_{P\in\mathcal E_t}\sigma(P).
\end{equation*}
Since~$V(A_u)=\sum_{i\in A_u}V_i$, the patch intervals also split the
Feynman--Kac integral:
\begin{equation*}
\int_0^tV(A_u)\,du
=
\sum_{P\in\mathcal B_t}
V_{i(P)}\int_{s(P)}^{e(P)}X_u^P\,du
+
\sum_{P\in\mathcal E_t}
V_{i(P)}\int_{s(P)}^tX_u^P\,du.
\end{equation*}
The end patches record the active set at time~$t$, so
\begin{equation*}
\chi_{A_t}(\eta)
=
\prod_{P\in\mathcal E_t}
\eta\bb{i(P)}^{X_t^P}.
\end{equation*}
Combining these identities gives the pathwise factorization
\begin{equation}
\label{eq:pathwise-patch-factorization}
\sigma_t
\exp\bb{\int_0^tV(A_u)\,du}
\chi_{A_t}(\eta)
=
\prod_{P\in\mathcal B_t}F(P)
\prod_{P\in\mathcal E_t}F\bb{\eta\bb{i(P)},P}.
\end{equation}
Condition~\eqref{eq:fk-integrability} makes the left-hand side of
\eqref{eq:pathwise-patch-factorization} integrable. Conditioning on~$\mathcal
G_t$ and applying Theorem~\ref{thm:patch-factorization} therefore replaces
each local factor by its consistent patch expectation. Taking expectation and
using Theorem~\ref{thm:monomial-duality}
proves~\eqref{eq:patch-representation}.
\end{proof}

\paragraph{Product initial laws.}
\label{subsec:general-initial-laws}

Integrating~\eqref{eq:patch-representation} gives the second assertion of
Theorem~\ref{thm:patch-representation}. For the Bernoulli product
measure~$\mu_{\mb p}$ with calm-state density
profile~$\mb p=(p_i)_{i\in\Lambda}$, the end-patch sites are distinct and the
end contributions are affine, so
\begin{equation}
\label{eq:patch-representation-product-law}
(\mu_{\mb p}P_t)(\chi_A)
=
\E_A\Cb{
\prod_{P\in\mathcal B_t}C(P)
\prod_{P\in\mathcal E_t}C\bb{p_{i(P)},P}
}.
\end{equation}
The skeleton retains the nonempty-target interactions that act on the dual
process. The other marked interactions between consecutive incoming or
outgoing successful interactions are integrated inside~$C(P)$ or~$C(z,P)$
before absolute values are taken. Thus~\eqref{eq:patch-representation} groups
the signed dual trajectory without discarding cancellations between successive
updates at a site.

\FloatBarrier
